% Source coverage: physical PDF pp. 376--381 (printed pp. 353--358), beginning at the Problems heading and continuing through the complete Chapter 21 reference list.
% Problems: 6. References: 57 entries (1--47 plus 3a, 5a, 20a, 20b, 27a--27c, 36a, 43a, and 44a). Unnumbered displays, footnotes, figures, tables, and centered dividers: none.
% Uncertainty: none. Visible source anomalies are preserved and recorded in the wave report.

\chapterbackmatter{Problems}

\begin{enumerate}
\item Calculate the effective ghost Lagrangian in a `generalized unitarity
gauge,' with $B[f]$ given as usual by
Eq.~\hyperref[eq:15.5.22]{(15.5.22)}, but now where
$f_\alpha(x)=i\phi_n(x)(t_\alpha)_{nm}
\langle\phi_m(0)\rangle_{\rm VAC}$, with $\phi_n$ real scalar fields, and
$t_\alpha$ imaginary antisymmetric matrices representing the Lie algebra of
the gauge group. What is the ghost propagator? Is this part of the
Lagrangian renormalizable?

\item What would be the effect in the $SU(2)\times U(1)$ electroweak theory if
the gauge symmetry were broken by the vacuum expectation value of a field
$\phi_3$ belonging to a real triplet
$\vec\phi=(\phi^+,\phi^0,\phi^-)$ instead of the usual complex doublet
$(\phi^0,\phi^-)$?

\item Consider the standard electroweak theory, with a single scalar doublet.
In one-loop order, calculate the effect of $Z^0$ and neutral scalar exchange
on the anomalous magnetic moment of the muon.

\item What is the lowest-order magnetic moment of the $W^+$ and $Z^0$
particles in the standard electroweak theory?

\item What would be the effect of the discovery of a fourth generation of
quarks and leptons on the predictions of the unified theories of strong and
electroweak interactions discussed in
\hyperref[sec:21.5]{Section 21.5}?

\item Suppose that several fields with incommensurate values of the electric
charge had nonvanishing vacuum expectation values in a superconductor. What
effect would this have on the properties of superconductors discussed in
\hyperref[sec:21.6]{Section 21.6}?
\end{enumerate}

\chapterbackmatter{References}

\begin{enumerate}
\item \phantomsection\label{ch21-ref-1}
J. Goldstone, \emph{Nuovo Cimento} \textbf{9}, 154 (1961); J. Goldstone,
A. Salam, and S. Weinberg, \emph{Phys. Rev.} \textbf{127}, 965 (1962).

\item \phantomsection\label{ch21-ref-2}
P. W. Higgs, \emph{Phys. Lett.} \textbf{12}, 132 (1964);
\emph{Phys. Rev. Lett.} \textbf{13}, 508 (1964);
\emph{Phys. Rev.} \textbf{145}, 1156 (1966); F. Englert and R. Brout,
\emph{Phys. Rev. Lett.} \textbf{13}, 321 (1964); G. S. Guralnik,
C. R. Hagen, and T. W. B. Kibble, \emph{Phys. Rev. Lett.} \textbf{13},
585 (1964); T. W. B. Kibble, \emph{Phys. Rev.} \textbf{155}, 1554
(1967); S. Weinberg, \emph{Phys. Rev. Lett.} \textbf{18}, 507 (1967):
footnote 7.

\item \phantomsection\label{ch21-ref-3}
S. Weinberg, \emph{Phys. Rev. Lett.} \textbf{19}, 1264 (1967). A. Salam,
in \emph{Elementary Particle Physics}, ed. N. Svartholm (Almqvist and
Wiksells, Stockholm, 1968): p. 367.

\item[3a.] \phantomsection\label{ch21-ref-3a}
The Feynman rules for general gauge theories in unitarity gauge were given
by S. Weinberg, \emph{Phys. Rev.} \textbf{D7}, 1068 (1973).

\item \phantomsection\label{ch21-ref-4}
G. 't Hooft, \emph{Nucl. Phys.} \textbf{B35}, 167 (1971); B. W. Lee,
\emph{Phys. Rev.} \textbf{D5}, 823 (1972).

\item \phantomsection\label{ch21-ref-5}
K. Fujikawa, B. W. Lee, and A. Sanda, \emph{Phys. Rev.} \textbf{D6},
2923 (1972).

\item[5a.] \phantomsection\label{ch21-ref-5a}
B. W. Lee and J. Zinn-Justin, \emph{Phys. Rev.} \textbf{D5}, 3121, 3137
(1972); \emph{Phys. Rev.} \textbf{D7}, 1049 (1972); G. 't Hooft and
M. Veltman, \emph{Nucl. Phys.} \textbf{B50}, 318 (1972); B. W. Lee,
\emph{Phys. Rev.} \textbf{D9}, 933 (1974).

\item \phantomsection\label{ch21-ref-6}
H. Georgi and S. L. Glashow, \emph{Phys. Rev. Lett.} \textbf{28}, 1494
(1972).

\item \phantomsection\label{ch21-ref-7}
T. D. Lee and C. N. Yang, \emph{Phys. Rev.} \textbf{98}, 101 (1955).

\item \phantomsection\label{ch21-ref-8}
Models with an incomplete $SU(2)\times U(1)$ symmetry were proposed by
S. L. Glashow, \emph{Nucl. Phys.} \textbf{22}, 579 (1961); A. Salam and
J. Ward, \emph{Phys. Lett.} \textbf{13}, 168 (1964). Also see
J. Schwinger, \emph{Ann. Phys. (N.Y.)} \textbf{2}, 407 (1957).

\item \phantomsection\label{ch21-ref-9}
Speculations about possible neutral currents go back to work of G. Gamow
and E. Teller, \emph{Phys. Rev.} \textbf{51}, 288 (1937); N. Kemmer,
\emph{Phys. Rev.} \textbf{52}, 906 (1937); G. Wentzel,
\emph{Helv. Phys. Acta} \textbf{10}, 108 (1937); S. Bludman,
\emph{Nuovo Cimento} \textbf{9}, 433 (1958); J. Leites-Lopes,
\emph{Nucl. Phys.} \textbf{8}, 234 (1958).

\item \phantomsection\label{ch21-ref-10}
Neutral currents were first observed in a single $\bar\nu_\mu e^-$
scattering event at CERN, by F. J. Hasert \emph{et al.}
(Aachen--Brussels--CERN--Ecole Polytechnique--Milan--Orsay--London
collaboration), \emph{Phys. Lett.} \textbf{46}, 121 (1973).

\item \phantomsection\label{ch21-ref-11}
N. Cabibbo, \emph{Phys. Rev. Lett.} \textbf{10}, 531 (1963).

\item \phantomsection\label{ch21-ref-12}
I. S. Towner, E. Hagberg, J. C. Hardy, V. T. Koslowsky, and G. Savard,
Chalk River preprint nucl-th/9507005 (1995).

\item \phantomsection\label{ch21-ref-13}
S. L. Glashow, J. Iliopoulos, and L. Maiani, \emph{Phys. Rev.}
\textbf{D2}, 1285 (1970).

\item \phantomsection\label{ch21-ref-14}
S. Weinberg, \emph{Phys. Rev. Lett.} \textbf{27}, 1688 (1971);
\emph{Phys. Rev.} \textbf{5}, 1413 (1972).

\item \phantomsection\label{ch21-ref-15}
J. J. Aubert \emph{et al.}, \emph{Phys. Rev. Lett.} \textbf{33}, 1404
(1974); J. E. Augustin \emph{et al.}, \emph{Phys. Rev. Lett.}
\textbf{33}, 1406 (1974).

\item \phantomsection\label{ch21-ref-16}
M. L. Perl \emph{et al.}, \emph{Phys. Rev. Lett.} \textbf{35}, 1489
(1975).

\item \phantomsection\label{ch21-ref-17}
S. W. Herb \emph{et al.}, \emph{Phys. Rev. Lett.} \textbf{39}, 252
(1977).

\item \phantomsection\label{ch21-ref-18}
F. Abe \emph{et al.}, \emph{Phys. Rev. Lett.} \textbf{74}, 2626 (1995);
S. Abachi \emph{et al.}, \emph{Phys. Rev. Lett.} \textbf{74}, 2632
(1995).

\item \phantomsection\label{ch21-ref-19}
The values given by the CDF and D0 collaborations at Fermilab are
respectively $176\pm8\pm10$ GeV (F. Abe \emph{et al.},
\hyperref[ch21-ref-18]{Ref.~18}.) and
$199^{+19}_{-21}\pm22$ GeV (S. Abachi \emph{et al.},
\hyperref[ch21-ref-18]{Ref.~18}), with the first and second errors
statistical and systematic, respectively. These results are interpreted to
give a combined value of $181\pm12$ GeV by J. Ellis, G. L. Fogli, and
E. Lisi, CERN--BARI preprint hep-ph/9507424.

\item \phantomsection\label{ch21-ref-20}
M. Kobayashi and K. Maskawa, \emph{Prog. Theor. Phys.} \textbf{49}, 282
(1972).

\item[20a.] \phantomsection\label{ch21-ref-20a}
F. J. Gilman, K. Kleinknecht, and B. Renk, Carnegie Mellon--Mainz
preprint CMU-HUP95-19-DOE-ER/40682-107 (1995), to be published in the
1995 \emph{Review of Particle Properties}.

\item[20b.] \phantomsection\label{ch21-ref-20b}
S. Weinberg, \emph{Phys. Rev. Lett.} \textbf{37}, 657 (1976). For earlier
models in which scalar fields are responsible for CP and T
non-conservation see T. D. Lee, \emph{Phys. Rev.} \textbf{D8}, 1226
(1973); \emph{Phys. Rep.} \textbf{9C}, 143 (1974).

\item \phantomsection\label{ch21-ref-21}
F. J. Hasert \emph{et al.}, \emph{Phys. Lett.} \textbf{46B}, 138 (1973);
P. Musset, \emph{Jour. de Physique} \textbf{11/12}, T34 (1973). Neutral
current events were seen at about the same time by the
Harvard--Pennsylvania--Wisconsin--Fermilab group at Fermilab, but
publication of their paper was delayed, so they took the opportunity to
rebuild their detector, and at first did not find the same signal. Evidence
for neutral currents was published by this group in A. Benvenuti
\emph{et al.}, \emph{Phys. Rev. Lett.} \textbf{32}, 800 (1974).

\item \phantomsection\label{ch21-ref-22}
G. Arnison \emph{et al.} \emph{Phys. Lett.} \textbf{122B}, 103 (1983);
\textbf{126B}, 398 (1983); \textbf{129B}, 273 (1983); \textbf{134B}, 469
(1984); \textbf{147B}, 241 (1984).

\item \phantomsection\label{ch21-ref-23}
F. Abe \emph{et al.} (CDF collaboration), \emph{Phys. Rev. Lett.}
\textbf{75}, 11 (1995). The $W$ mass is measured here in observations of
the decays $W\to\mu+\nu$ and $W\to e+\nu$.

\item \phantomsection\label{ch21-ref-24}
CERN report LEPEWWG/95-01, unpublished (1995).

\item \phantomsection\label{ch21-ref-25}
M. Green and M. Veltman, \emph{Nucl. Phys.} \textbf{B169}, 137 1980);
V. A. Novikov, L. B. Okun', and M. I. Vysotsky, \emph{Nucl. Phys.}
\textbf{B397}, 35 (1992). For later studies, see P. Langacker, in
\emph{Precision Tests of the Standard Model} (World Scientific, Singapore,
1994); P. Bamert, C. P. Burgess, and I. Maksymyk,
McGill--Neuchatal--Texas preprint hep-ph/9505339 (1995); W. Hollik,
Karlsruhe preprint hep-ph/9507406 (1995).

\item \phantomsection\label{ch21-ref-26}
S. Fanchotti, B. Kniehl, and A. Sirlin, \emph{Phys. Rev.} \textbf{D48},
307 (1973), and references cited therein.

\item \phantomsection\label{ch21-ref-27}
J. Ellis, G. L. Fogli, and E. Lisi, \hyperref[ch21-ref-19]{Ref.~19}.

\item[27a.] \phantomsection\label{ch21-ref-27a}
S. Weinberg, \emph{Phys. Rev. Lett.} \textbf{43}, 1566 (1979).

\item[27b.] \phantomsection\label{ch21-ref-27b}
For instance, in the so-called see-saw mechanism, a neutrino mass of this
order would be produced by exchange of a heavy neutral lepton of mass $M$;
see M. Gell-Mann, P. Ramond, and R. Slansky, in \emph{Supergravity},
ed. P. van Nieuwenhuizen and D. Freedman (North Holland, Amsterdam, 1979):
p. 315; T. Yanagida, \emph{Prog. Theor. Phys.} \textbf{B135}, 66 (1978).

\item[27c.] \phantomsection\label{ch21-ref-27c}
S. Weinberg, \hyperref[ch21-ref-27a]{Ref.~27a}; F. Wilczek and A. Zee,
\emph{Phys. Rev. Lett.} \textbf{43}, 1571 (1979). Also see S. Weinberg,
\emph{Phys. Rev.} \textbf{D22}, 1694 (1980).

\item \phantomsection\label{ch21-ref-28}
This general formalism was described by S. Weinberg, \emph{Phys. Rev.}
\textbf{D13}, 974 (1976). For earlier work on dynamical symmetry breaking
in gauge theories, see R. Jackiw and K. Johnson, \emph{Phys. Rev.}
\textbf{D8}, 2386 (1973); J. M. Cornwall and R. E. Norton,
\emph{Phys. Rev.} \textbf{D8}, 3338 (1973).

\item \phantomsection\label{ch21-ref-29}
This interpretation of the results of the standard model is due to
L. Susskind, \hyperref[ch21-ref-30]{Ref.~30}.

\item \phantomsection\label{ch21-ref-30}
S. Weinberg, \hyperref[ch21-ref-28]{ref.~28} and \emph{Phys. Rev.}
\textbf{D19}, 1277 (1979); L. Susskind, \emph{Phys. Rev.} \textbf{D19},
2619 (1979). The term `technicolor' is due to Susskind.

\item \phantomsection\label{ch21-ref-31}
S. Dimopoulos and L. Susskind, \emph{Nucl. Phys.} \textbf{B15}, 237
(1979); E. Eichten and K. Lane, \emph{Phys. Lett.} \textbf{90B}, 125
(1980); S. Weinberg, unpublished work quoted by Eichten and Lane. For a
review, see E. Farhi and L. Susskind, \emph{Phys. Rep.} \textbf{74}, 277
(1981).

\item \phantomsection\label{ch21-ref-32}
S. Weinberg, \emph{Phys. Rev.} \textbf{D5}, 1962 (1972).

\item \phantomsection\label{ch21-ref-33}
J. Pati and A. Salam, \emph{Phys. Rev. Lett.} \textbf{31}, 275 (1973).

\item \phantomsection\label{ch21-ref-34}
H. Georgi and S. L. Glashow, \emph{Phys. Rev. Lett.} \textbf{32}, 438
(1974).

\item \phantomsection\label{ch21-ref-35}
H. Georgi, in \emph{Particles and Fields -- 1974}, ed. C. Carlson
(Amer. Inst. of Physics, New York, 1975).

\item \phantomsection\label{ch21-ref-36}
H. Georgi, H. R. Quinn, and S. Weinberg, \emph{Phys. Rev. Lett.}
\textbf{33}, 451 (1974).

\item[36a.] \phantomsection\label{ch21-ref-36a}
S. Dimopoulos and H. Georgi, \emph{Nucl. Phys.} \textbf{B193}, 150
(1981); J. Ellis, S. Kelley, and D. V. Nanopoulos, \emph{Phys. Lett.}
\textbf{B260}, 131 (1991); U. Amaldi, W. de Boer, and H. Furstmann,
\emph{Phys. Lett.} \textbf{B260}, 447 (1991). For more recent analyses of
the data, see P. Langacker and N. Polonsky, \emph{Phys. Rev.} \textbf{D47},
4028 (1993); \textbf{D49}, 1454 (1994); L. J. Hall and U. Sarid,
\emph{Phys. Rev. Lett.} \textbf{70}, 2673 (1993).

\item \phantomsection\label{ch21-ref-37}
J. Bardeen, L. N. Cooper, and J. R. Schrieffer, \emph{Phys. Rev.}
\textbf{108}, 1175 (1957).

\item \phantomsection\label{ch21-ref-38}
P. M. Anderson \emph{Phys. Rev.} \textbf{130}, 439 (1963).

\item \phantomsection\label{ch21-ref-39}
S. Weinberg, \emph{Prog. Theor. Phys. Suppl.} No. 86, 43 (1986).

\item \phantomsection\label{ch21-ref-40}
B. D. Josephson, \emph{Phys. Lett.} \textbf{1}, 25 (1962).

\item \phantomsection\label{ch21-ref-41}
V. L. Ginzburg and L. D. Landau, \emph{JETP (USSR)} \textbf{20}, 1064
(1950).

\item \phantomsection\label{ch21-ref-42}
L. P. Gor'kov, \emph{Soviet Phys. JETP} \textbf{9}, 1364 (1959).

\item \phantomsection\label{ch21-ref-43}
A. A. Abrikosov, \emph{Soviet Phys. JETP} \textbf{5}, 1174 (1957).

\item[43a.] \phantomsection\label{ch21-ref-43a}
E. B. Bogomol'nyi, \emph{Sov. J. Nucl. Phys.} \textbf{24}, 449 (1976).
For numerical calculations, see E. B. Bogomol'nyi and A. I. Vainshtein,
\emph{Sov. J. Nucl. Phys.} \textbf{23}, 588 (1976).

\item \phantomsection\label{ch21-ref-44}
G. Benfatto and G. Gallavotti, \emph{J. Stat. Phys.} \textbf{59}, 541
(1990); \emph{Phys. Rev.} \textbf{42}, 9967 (1990); J. Feldman and
E. Trubowitz, \emph{Helv. Phys. Acta} \textbf{63}, 157 (1990);
\textbf{64}, 213 (1991); \textbf{65}, 679 (1992); R. Shankar,
\emph{Physica} \textbf{A177}, 530 (1991); \emph{Rev. Mod. Phys.}
\textbf{66}, 129 (1993); J. Polchinski, in \emph{Recent Developments in
Particle Theory, Proceedings of the 1992 TASI}, eds. J. Harvey and
J. Polchinski (World Scientific, Singapore, 1993).

\item[44a.] \phantomsection\label{ch21-ref-44a}
S. Weinberg, \emph{Nucl. Phys.} \textbf{B413}, 567 (1994).

\item \phantomsection\label{ch21-ref-45}
R. L. Stratonovich, \emph{Sov. Phys. Dokl.} \textbf{2}, 416 (1957);
J. Hubbard, \emph{Phys. Rev. Lett.} \textbf{3}, 77 (1959).

\item \phantomsection\label{ch21-ref-46}
A. A. Abrikosov, L. P. Gor'kov, I. Ye. Dzyaloshinskii,
\emph{Quantum Field Theoretical Methods in Statistical Physics}, 2nd edn.,
transl. by D. E. Brown (Pergamon Press, Oxford, 1965).

\item \phantomsection\label{ch21-ref-47}
In order to give a precise meaning to the mass $M$ it is necessary to take
into account both two-loop corrections to the renormalization group
equations \hyperref[eq:21.5.6]{(21.5.6)}--\hyperref[eq:21.5.8]{(21.5.8)},
and one-loop `threshold corrections' to condition
\hyperref[eq:21.5.5]{(21.5.5)}. See S. Weinberg,
\emph{Phys. Lett.} \textbf{82B}, 387 (1979).
\end{enumerate}
